% ─── Chapter 3: Methods ────────────────────────────────────────────
\chapter{Methods}
\label{ch:methods}

This chapter describes the reinforcement learning framework, model architecture,
training infrastructure, and the iterative development methodology used to
diagnose and fix issues in the training pipeline.

\section{Reinforcement Learning Framework}

\subsection{MDP Formulation}

The Pok\'{e}mon battle environment is modelled as a Markov Decision Process
(MDP) with the following components:

\begin{center}
\small
\begin{tabularx}{\textwidth}{l X}
  \toprule
  \textbf{Component} & \textbf{Definition} \\
  \midrule
  State ($s$)       & Dictionary observation: token matrix $(13 \times 164)$,
                      species/item/ability ID vectors $(13,)$ each, action
                      mask $(14,)$.  Full specification in
                      Section~\ref{sec:architecture}. \\
  Action ($a$)       & Discrete choice from 14 possible actions (4 moves,
                      4 gimmick-enhanced moves, 6 switches), filtered by the
                      action mask to only legal options. \\
  Reward ($r$)       & Shaped reward computed each step from HP changes,
                      fainting events, and a terminal outcome bonus/penalty.
                      Detailed in Section~\ref{sec:reward_design}. \\
  Transition         & Deterministic given both players' actions;
                      stochasticity arises from the opponent's policy and
                      damage variance in the simulator. \\
  Discount ($\gamma$) & 0.97 based on a hyper parameter sweep. \\
  \bottomrule
\end{tabularx}
\end{center}

An episode corresponds to a single battle.  It terminates when one side has no
remaining conscious Pok\'{e}mon (win or loss).  There is no explicit time limit
or truncation.

\subsection{PPO Algorithm}

The agent is trained with \textbf{Proximal Policy Optimization} (PPO), an
on-policy actor-critic algorithm that balances sample efficiency with training
stability through a clipped surrogate objective.  PPO was chosen for the
following reasons:

\begin{itemize}
  \item \textbf{Stability:} The clipped objective prevents destructively large
        policy updates, which is critical when learning in a sparse-reward,
        high-variance environment.
  \item \textbf{Distributed rollouts:} PPO naturally integrates with Ray
        RLlib's distributed rollout architecture, enabling many parallel
        battles.
  \item \textbf{Action masking:} PPO's categorical distribution can be
        combined with action masking (setting invalid-action logits to a large
        negative value) without modifying the algorithm itself.
\end{itemize}

The policy network (actor) outputs a categorical distribution over the 14
actions, and the value network (critic) outputs a scalar state-value estimate
$V(s)$.  Both share a common transformer trunk (see
Section~\ref{sec:architecture}).

\subsection{Reward Design}
\label{sec:reward_design}

The reward function combines several components to provide dense learning signal
while preserving the primacy of winning.  All rewards are multiplied by a
global \texttt{reward\_scale}$\,=0.1$ or 0.05 before being returned to the agent, which
scales returns from approximately $[-15, +15]$ to $[-1.5, +1.5]$.  The
rationale for this scaling is discussed in
Section~\ref{sec:reward_scaling}.

\begin{center}
\small
\begin{tabularx}{\textwidth}{l r X}
  \toprule
  \textbf{Component} & \textbf{Default} & \textbf{Description} \\
  \midrule
  Victory reward     & $+10.0$  & Agent wins the battle (terminal step). \\
  Defeat penalty     & $-10.0$  & Agent loses the battle (terminal step). \\
  HP differential    & $\times 2.0$ &
    $(\text{our\_team\_HP} - \text{opp\_team\_HP}) \times w$ as a shaped
    step reward.  Positive when the agent is ahead, negative when behind. \\
  Opponent fainted   & $+5.0$   & Per opponent Pok\'{e}mon knocked out. \\
  Own fainted        & $-5.0$   & Per agent Pok\'{e}mon that faints. \\
  Step penalty       & $-0.005$ & Per step, encourages battle efficiency
                                    (fewer turns). \\
  Matchup quality    & $\times 0.2$ &
    $+1.0$ if a super-effective move is available, $-0.5$ if only resisted
    moves, $-1.0$ if immune.  Rewards good type-readiness. \\
  Action quality     & $\times 0.3$ &
    Penalty for clearly suboptimal moves; bonus for resisting the opponent's
    best move. \\
  \bottomrule
\end{tabularx}
\end{center}

The reward is computed each step by \texttt{poke-env}'s
\texttt{reward\_computing\_helper}, parameterised by the active
\texttt{RewardConfig}.  At the terminal step, explicit outcome rewards are
applied on top.  Reward parameters are adjusted across curriculum stages (see
below). These were gradually introduced during development in descending order.

\subsection{Curriculum Learning}
\label{sec:curriculum}

Due to the complexity of the full challenge, training follows a \textbf{curriculum
learning} approach where the agent is first exposed to easy opponents and
gradually faces stronger ones.  The curriculum currently has four stages:

\begin{enumerate}
  \item \textbf{Moves \& Switches Warmup.}
        Opponent mix: 40\,\% random, 30\,\% \texttt{random\_no\_switch},
        30\,\% self-play.  Promotion threshold: 65\,\% win rate.
        Asymmetric rewards (victory $+8$, defeat $-10$) encourage caution.
        Higher HP weight (3.0) and faint values ($\pm 6.0$) provide strong
        per-step signal for learning basic move selection.

  \item \textbf{Random, More Moves.}
        Opponent mix: 70\,\% self-play, 10\,\% random,
        20\,\% \texttt{random\_no\_switch}.  Promotion threshold: 20\,\% win
        rate (intentionally low: the goal is exposure to self-play, not
        mastery).  Victory reward raised to $+10$.  Large minimum sample
        requirement (5\,000 episodes) ensures sufficient exposure before
        progressing.

  \item \textbf{Self-Play.}
        Opponent mix: 60\,\% \texttt{random\_no\_switch},
        40\,\% random.  Promotion threshold: 70\,\% win rate.
        After the self-play exposure stage, the agent trains against
        non-switching random opponents to solidify move selection skills.
        Action quality weight increased to 0.4 to refine move choice.

  \item \textbf{Mixed Final.}
        Opponent mix: 50\,\% heuristic, 30\,\% self-play,
        20\,\% \texttt{random\_no\_switch}.  Terminal stage (promotion
        threshold 101\,\%, i.e.\ never promotes).  The agent faces the
        strongest opponent mix, including the heuristic AI which considers
        type effectiveness and switches strategically.
\end{enumerate}

% TODO: If a fifth stage is added, update the list above.

\subsubsection{Per-Stage Reward Variants}

Reward parameters are adjusted across stages to reflect changing difficulty:

\begin{center}
\small
\begin{tabularx}{\textwidth}{l c c c c}
  \toprule
  \textbf{Parameter} & \textbf{Stage 1} & \textbf{Stage 2} & \textbf{Stage 3} & \textbf{Stage 4} \\
  \midrule
  Victory / Defeat   & $+8$ / $-10$ & $\pm 10$ & $\pm 10$ & $\pm 10$ \\
  HP weight          & 3.0   & 3.0   & 3.0   & 3.0   \\
  Opp.\ fainted      & $+6.0$ & $+6.0$ & $+6.0$ & $+6.0$ \\
  Own fainted        & $-6.0$ & $-6.0$ & $-6.0$ & $-6.0$ \\
  Step penalty       & $-0.01$ & $-0.01$ & $-0.01$ & $-0.01$ \\
  Matchup weight     & 0.1   & 0.1   & 0.1   & 0.1   \\
  Action quality wt. & 0.3   & 0.3   & 0.4   & 0.4   \\
  \bottomrule
\end{tabularx}
\end{center}

Note that Stage~1 uses asymmetric rewards (victory $+8$ vs.\ defeat $-10$) to
discourage early aggression, and a higher step penalty ($-0.01$) to encourage
efficient play from the start.

Promotion decisions use a rolling window of the last 300~episodes (minimum 50
samples required before promotion is considered; minimum 300 episodes must
elapse before any promotion). Total possible rewards is by design made in a way that can only increase with difficulty.

\subsection{Self-Play Mechanism}
\label{sec:self_play}

Self-play is implemented through a \texttt{SelfPlayPlayer} class that:

\begin{enumerate}
  \item Loads the latest model weights from
        \texttt{checkpoints/selfplay\_latest.pt} on first use.
  \item Re-checks the weights file each turn, enabling automatic updates
        when new checkpoints are saved.
  \item Maintains per-battle LSTM state for temporal memory.
  \item Converts the agent's compressed 14-action output to native
        \texttt{poke-env} \texttt{BattleOrder} objects.
  \item Falls back to \texttt{RandomPlayer} behaviour on inference errors.
\end{enumerate}

At each checkpoint, the training pipeline exports the current model weights to
the self-play weights path, making them immediately available to all
\texttt{SelfPlayPlayer} instances across workers.  This creates a natural
curriculum where the agent continuously adapts to an opponent of increasing
skill.

% TODO: Add talking point about self-play training results once available.
% Specifically: win rate curves for self-play stage, comparison with
% non-self-play training, and any observed oscillation or instability.

\subsection{Opponent Types}

The environment supports four opponent categories:

\begin{center}
\small
\begin{tabularx}{\textwidth}{l l X}
  \toprule
  \textbf{Key} & \textbf{Player Class} & \textbf{Behaviour} \\
  \midrule
  \texttt{random\_no\_switch} & \texttt{RandomNoSwitchPlayer} &
    Uniform random among legal \emph{move} orders; switches only when forced
    (no moves available or forced by simulator). \\
  \texttt{random} & \texttt{RandomPlayer} &
    Uniform random over all legal orders (moves \emph{and} switches). \\
  \texttt{heuristic} & \texttt{SimpleHeuristicsPlayer} &
    Rule-based AI: considers type effectiveness, switches when at a
    disadvantage. \\
  \texttt{self} & \texttt{SelfPlayPlayer} &
    Uses the training model's weights for inference with hot-reload support. \\
  \bottomrule
\end{tabularx}
\end{center}

In the embedding layer, \texttt{random\_no\_switch} is mapped to the same
bucket as \texttt{random} for the opponent-type feature, avoiding distribution
shift at the observation level when mixing or swapping baselines. From best to worst in terms of performance: heuristic, random no switch and random.

\section{Model Architecture}
\label{sec:architecture}

\subsection{Overview}

The system consists of three major components: (1)~a \textbf{custom Gymnasium
environment} wrapping \texttt{poke-env} and Pok\'{e}mon Showdown, (2)~a
\textbf{transformer-based neural network} that processes battle observations
and outputs action logits plus a state-value estimate, and (3)~a
\textbf{distributed training pipeline} built on Ray RLlib.  This section
details the first two; the training infrastructure is covered in
Section~\ref{sec:infrastructure}.

\subsection{Observation Space}

The observation is a dictionary with five keys:

\begin{center}
\small
\begin{tabularx}{\textwidth}{l l X}
  \toprule
  \textbf{Key} & \textbf{Shape / Dtype} & \textbf{Content} \\
  \midrule
  \texttt{obs}          & $(13,\,164)$ float32 &
    Token matrix: 1~global token + 6~own-team tokens + 6~opponent-team
    tokens.  Each token encodes presence flags, HP, base stats, types,
    status, moves, and more (see Appendix~\ref{ch:appendix_obs}). \\
  \texttt{species}      & $(13,)$ int32 &
    Species ID per token (dictionary lookup, vocabulary size 1\,151). \\
  \texttt{items}        & $(13,)$ int32 &
    Held-item ID per token (vocabulary size 1\,110). \\
  \texttt{abilities}    & $(13,)$ int32 &
    Ability ID per token (vocabulary size 216). \\
  \texttt{action\_mask} & $(14,)$ float32 &
    Binary mask: 1.0 for valid actions, 0.0 for invalid. \\
  \bottomrule
\end{tabularx}
\end{center}

\subsubsection{Token Layout}

The 13 tokens are organised as follows:

\begin{center}
\small
\begin{tabularx}{0.7\textwidth}{c l l}
  \toprule
  \textbf{Index} & \textbf{Role} & \textbf{Content} \\
  \midrule
  0        & Global      & Weather, terrain, side conditions, battle extras \\
  1        & Our active  & Active Pok\'{e}mon on the agent's side \\
  2--6     & Our bench   & Remaining team Pok\'{e}mon \\
  7        & Opp.\ active & Active opponent Pok\'{e}mon \\
  8--12    & Opp.\ bench  & Remaining opponent team \\
  \bottomrule
\end{tabularx}
\end{center}

Each Pok\'{e}mon token is 164~dimensions.  The full per-token feature
breakdown is provided in Appendix~\ref{ch:appendix_obs}.

\subsection{Embedding Layers}

High-cardinality categorical entities are mapped to learned embeddings:

\begin{center}
\small
\begin{tabularx}{0.7\textwidth}{l r r}
  \toprule
  \textbf{Entity} & \textbf{Vocabulary Size} & \textbf{Embedding Dim} \\
  \midrule
  Species   & 1\,151 & 32 \\
  Items     & 1\,110 & 32 \\
  Abilities & 216    & 32 \\
  \bottomrule
\end{tabularx}
\end{center}

Padding index~0 represents absent entities (e.g.\ a bench slot with no
Pok\'{e}mon).  Each token's base 164-dim vector is concatenated with three
32-dim embedding vectors ($3 \times 32 = 96$), yielding 260~dimensions total.
A linear projection followed by LayerNorm maps this to the transformer's hidden
dimension (512 by default).

\subsubsection{From Hashes to Dictionaries}

Earlier versions used Adler-32 hashing with modulo 20\,000 to compress entity
strings into integer IDs.  This caused approximately 7.5\,\% collision rate
among 1\,800 unique strings.  The current system uses stable dictionary
lookups generated from Showdown's data files, eliminating collisions entirely
and ensuring that each entity receives a unique embedding.

\subsection{Position \& Role Embeddings}

The transformer processes a set of 13 tokens.  Without positional information,
self-attention is permutation-invariant, meaning the model cannot distinguish between
active and bench Pok\'{e}mon.  Two sets of learnable embeddings address this:

\begin{itemize}
  \item \textbf{Position embeddings:} One learned vector per token position
        (13~positions $\to$ hidden dim).  Added after input projection.
  \item \textbf{Role embeddings:} Five coarse roles: CLS (global), our
        active, our bench, opponent active, opponent bench.  Shared across
        positions within the same role.
\end{itemize}

Both are added to the projected token vectors before the transformer encoder.

\subsection{Transformer Encoder}

The core of the policy network is a standard \texttt{nn.TransformerEncoder}
with the following configuration:

\begin{center}
\small
\begin{tabularx}{0.6\textwidth}{l r}
  \toprule
  \textbf{Parameter} & \textbf{Value} \\
  \midrule
  Hidden dimension     & 512 \\
  Attention heads      & 8 \\
  Transformer layers   & 1 \\
  Feedforward expansion & $\times 4$ (2\,048) \\
  Activation           & GELU \\
  Dropout              & 0.1 \\
  Normalisation        & Pre-norm (LayerNorm before attention/FFN) \\
  \bottomrule
\end{tabularx}
\end{center}

The decision to use a single transformer layer was driven by an attention
collapse analysis (see Section~\ref{sec:dev_methodology}), which showed that deeper
post-norm transformers suffered from gradient vanishing in layers 2--3,
effectively reducing model depth to two functional layers.

\subsubsection{Cross-Team Attention Bias}

A learnable \textbf{attention bias} is added to the self-attention scores to
encourage the model to attend to the opponent's active Pok\'{e}mon: a
critical source of tactical information:

\begin{center}
\small
\begin{tabularx}{0.7\textwidth}{l l r}
  \toprule
  \textbf{From} & \textbf{To} & \textbf{Bias} \\
  \midrule
  CLS (0)         & Opp.\ active (7)  & $+2.0$ \\
  Our active (1)  & Opp.\ active (7)  & $+2.0$ \\
  Opp.\ active (7) & Our active (1)   & $+1.0$ \\
  Opp.\ active (7) & Our bench (2--6) & $+0.5$ \\
  CLS (0)         & Opp.\ bench (8--12) & $+0.5$ \\
  \bottomrule
\end{tabularx}
\end{center}

This bias was introduced after analysis showed that the default transformer
progressively lost opponent attention in deeper layers.  With a single layer
and this bias, the model allocates 37--63\,\% of CLS attention to the
opponent's active Pok\'{e}mon.

% TODO: Insert MLflow screenshot or generated figure showing attention
% heatmap for the 1-layer model with cross-team bias.  This should show
% the attention distribution from CLS and our_active to all 13 tokens.

\subsection{Action Space}
\label{sec:action_space}

The agent operates in a \textbf{compressed action space} of 14~discrete
actions:

\begin{center}
\small
\begin{tabularx}{\textwidth}{c l l l}
  \toprule
  \textbf{Index} & \textbf{Category} & \textbf{Compressed} & \textbf{Native Mapping} \\
  \midrule
  0--3  & Regular moves & Use move 1--4    & Direct (or first legal variant) \\
  4--7  & Gimmick moves & Gimmick + move 1--4 & Auto-selects: Mega $\to$ Z $\to$ Dynamax \\
  8--13 & Switches      & Switch to slot 1--6 & Mapped to party slot index \\
  \bottomrule
\end{tabularx}
\end{center}

The native \texttt{poke-env} space has 22~actions (6~switches + 4~regular
moves + 12~gimmick variants).  Compression to 14 actions was motivated by:

\begin{itemize}
  \item \textbf{Reduced branching factor:} 14 vs.\ 22 simplifies exploration.
  \item \textbf{Logical grouping:} Moves first, then gimmicks, then switches
        provides a consistent structure the model can learn.
  \item \textbf{Gimmick auto-selection:} Rather than requiring the agent to
        learn separate Mega/Z/Dynamax variants, the compressed space presents
        a single ``gimmick'' action that resolves to whichever gimmick is
        legally available.
\end{itemize}

\subsubsection{Action Masking}

At each step, the environment provides an \texttt{action\_mask} of shape
$(14,)$ indicating which actions are legal.  The model applies a mask by
setting invalid-action logits to $-10^{8}$ before the softmax.

This ensures the agent never selects an illegal action, regardless of its
learned policy.  The mask accounts for fainted Pok\'{e}mon (cannot switch to
them), unavailable moves (no PP, disabled), and the absence of gimmick
conditions (e.g.\ Dynamax not available).

\subsection{LSTM Memory}

An optional \textbf{single-layer LSTM} (hidden size 512) can process the CLS
token sequence across turns within a battle, providing cross-turn temporal
memory:

\begin{itemize}
  \item \textbf{Input:} CLS token embedding from the transformer at each turn.
  \item \textbf{State:} Hidden state $h_t$ and cell state $c_t$ are carried
        across turns, reset at episode boundaries.
  \item \textbf{Purpose:} Enable the agent to remember previous turns (e.g.\
        which moves the opponent has used, whether Protect was used recently)
        rather than treating each turn independently.
\end{itemize}

The LSTM is enabled in the model code and processes both single-step and
sequence batches.  However, RLlib's connector path for recurrent state
propagation required additional engineering (see
Section~\ref{sec:lstm_memory_fix}).

\subsection{Policy \& Value Heads}

Both heads share the transformer trunk (and LSTM, if active).  The CLS token
output serves as the shared representation:

\begin{itemize}
  \item \textbf{Policy head:} $\text{hidden} \to \text{Linear}(512, 256) \to
        \text{GELU} \to \text{Linear}(256, 14)$.  Output is action logits,
        masked before softmax.
  \item \textbf{Value head:} $\text{hidden} \to \text{Linear}(512, 256) \to
        \text{GELU} \to \text{Linear}(256, 1)$.  Outputs scalar state-value
        estimate $V(s)$.
\end{itemize}

\subsection{Model Size}

The default (\texttt{standard}) configuration produces approximately
\textbf{10--15\,million parameters}.  Preset variations are detailed in
Section~\ref{sec:infrastructure}.

% TODO: Insert a figure showing the full model architecture as a block
% diagram: Input (obs dict) -> Embeddings -> Projection + Pos/Role ->
% Transformer (1L) -> [LSTM] -> Policy/Value heads.  Can be generated
% with draw.io, TikZ, or similar.

\section{Training Infrastructure}
\label{sec:infrastructure}

\subsection{Distributed Training Architecture}

Training is orchestrated by Ray RLlib, which distributes rollout collection
across multiple workers and performs SGD updates on a central learner:

\begin{center}
\small
\begin{tabularx}{\textwidth}{l X}
  \toprule
  \textbf{Component} & \textbf{Role} \\
  \midrule
  Main process   & Configures training, manages curriculum callbacks,
                   drives the training loop. \\
  GPU learner     & Performs PPO SGD updates on collected batches.
                   Uses PyTorch backend. \\
  Rollout workers & Each runs multiple environments, collects trajectories,
                   and sends batches to the learner. \\
  Showdown servers & 8 local Node.js instances on ports 8000--8007,
                     each hosting parallel battles. \\
  \bottomrule
\end{tabularx}
\end{center}

The default (\texttt{standard}) configuration uses 24~workers with 8
environments per worker, yielding \textbf{192~parallel battles}.  Each Showdown
server hosts approximately 24~concurrent battles.

\subsection{Hardware}

\begin{center}
\small
\begin{tabularx}{0.8\textwidth}{l l}
  \toprule
  \textbf{Resource} & \textbf{Specification} \\
  \midrule
  GPU         & NVIDIA RTX 5090 (\texttt{optimal} preset) \\
  RAM         & $\sim$44\,GB peak with 48 environments \\
  CPU         & Multi-core (workers run on CPU threads) \\
  Storage     & Local SSD for checkpoints and MLflow artifacts \\
  \bottomrule
\end{tabularx}
\end{center}

Memory management is critical: the project previously observed RAM growth from
8\,GB to 52\,GB over 5~hours due to stale entries in environment caches.  This
was mitigated with periodic cleanup and memory-leak monitoring.

\subsection{Configuration Presets}

Five hardware presets control model size, parallelism, and training duration:

\begin{center}
\small
\begin{tabularx}{\textwidth}{l c c c c}
  \toprule
  \textbf{Parameter} & \textbf{Quick} & \textbf{Standard} & \textbf{Memory Safe} & \textbf{Large} \\
  \midrule
  Hidden dim              & 128   & 512   & 256   & 768   \\
  Attention heads         & 8     & 8     & 4     & 12    \\
  Transformer layers      & 1     & 1     & 1     & 6     \\
  Embedding dim           & 32    & 32    & 32    & 32    \\
  LSTM hidden             & ---   & 512   & 256   & 768   \\
  Use LSTM                & No    & Yes   & No    & Yes   \\
  Train batch size        & 4\,096  & 6\,144  & 2\,048  & 16\,384 \\
  SGD minibatch size      & 256   & 512   & 128   & 512   \\
  Approx.\ params         & $\sim$1M & $\sim$10M & $\sim$4M & $\sim$42M \\
  \bottomrule
\end{tabularx}
\end{center}

The \texttt{quick} preset is used for smoke tests and debugging.  The
\texttt{standard} preset is the default for training runs.  The
\texttt{optimal} preset (not shown) matches \texttt{standard} but targets the
RTX~5090 specifically.

\subsection{PPO Hyperparameters}

Default hyperparameters across all presets (unless overridden):

\begin{center}
\small
\begin{tabularx}{0.7\textwidth}{l r}
  \toprule
  \textbf{Parameter} & \textbf{Value} \\
  \midrule
  Learning rate              & $2 \times 10^{-4}$ \\
  Discount ($\gamma$)        & 0.97 \\
  GAE lambda ($\lambda$)     & 0.87 \\
  Clip range ($\epsilon$)    & 0.08 \\
  Entropy coefficient        & 0.013 \\
  Value-function loss coeff. & 0.5 \\
  Value clip param           & 4.85 \\
  Gradient clipping          & 5.0 \\
  SGD iterations per batch   & 8 \\
  \bottomrule
\end{tabularx}
\end{center}

These are based on the first breakthrough run.
%https://mlflow-server.chozo.me/#/experiments/6/runs/84ec848a3d574816830d5367b4e360d6/model-metrics

\subsection{MLflow Experiment Tracking}

All training runs are logged to an \textbf{MLflow} tracking server under the
experiment name \texttt{Pokemon\_RL\_Battler}.  The following categories of
data are recorded:

\begin{itemize}
  \item \textbf{Configuration snapshot:} Full training config, model config,
        curriculum config, and hardware preset, stored as MLflow parameters.
  \item \textbf{Training metrics:} Policy loss, value loss, entropy,
        KL divergence, win rate, episode length, per-stage statistics, logged
        every training iteration.
  \item \textbf{Validation metrics:} Win rates against fixed opponents
        (random, heuristic), per-opponent-type breakdowns, fallback event
        counts, logged at scheduled evaluation intervals.
  \item \textbf{Diagnostic metrics:} Action-mask density, observation sample
        recordings, HP remaining distributions, faint counts per episode.
  \item \textbf{System metrics:} CPU/RAM/GPU utilisation, memory leak counters.
\end{itemize}

Metrics use nested keys (e.g.\
\texttt{validation/fixed\_paired/win\_rate\_vs\_heuristic}) to support
structured comparison in the MLflow UI.

% TODO: Insert MLflow screenshot showing the experiment comparison view
% with multiple runs.  Highlight win rate curves, loss curves, and
% the ability to compare configuration differences between runs.

\subsection{Checkpoint Management}

\begin{itemize}
  \item \textbf{Frequency:} Checkpoints saved every 200\,000 training steps.
  \item \textbf{Retention:} Up to 5~checkpoints retained (oldest automatically
        deleted).
  \item \textbf{Self-play export:} At each checkpoint, model weights are also
        exported to \texttt{checkpoints/selfplay\_latest.pt} for use by
        \texttt{SelfPlayPlayer} instances.
  \item \textbf{Resume:} Training can be resumed from the latest checkpoint
        using \texttt{-{}-resume-checkpoint latest}, preserving MLflow run
        continuity via \texttt{-{}-mlflow-run-id <RUN\_ID>}.
\end{itemize}

\subsection{Validation Pipeline}

Scheduled evaluation runs against fixed opponent configurations:

\begin{center}
\small
\begin{tabularx}{\textwidth}{l X}
  \toprule
  \textbf{Protocol} & \textbf{Description} \\
  \midrule
  Smoke test     & Quick sanity check: 10 episodes against random opponent. \\
  Fixed paired   & 20~frozen teams $\times$ 3~opponents (random + random no switch + heuristic) =
                   60~battles per evaluation run.  Primary comparison metric. \\
  Mirror         & Both sides use the same team; tests raw battling skill
                   without team advantage. \\
  \bottomrule
\end{tabularx}
\end{center}

Validation is triggered at regular intervals (every 100\,000 steps) during
training and can also be run manually against any checkpoint via
\texttt{scripts/validate\_checkpoint.py}.  Results are logged to MLflow with
per-opponent-type breakdowns.

\subsection{Training Throughput}
% Look into this! I don't know if its accurate

Observed throughput on the \texttt{standard} preset:
$\sim$800~environment samples per second (target: 1\,000+).  The primary
bottleneck is the cpu since all the servers are single core and run on cpu.

\section{Iterative Development Methodology}
\label{sec:dev_methodology}

The training pipeline was developed iteratively, with each phase addressing a
specific problem identified through diagnostic analysis.  The phases are
presented in chronological order.  Note that the curriculum configuration used
during early development phases (baseline through PPO pipeline fixes) differed
from the current four-stage system described in
Section~\ref{sec:curriculum}; the early curriculum used a simpler two-stage
setup with ``easy'' (random opponents only) and ``medium'' (50/50 random and
heuristic) stages.

\subsection{Baseline Establishment}

\textbf{Starting point.} The initial system used a 22-action native
\texttt{poke-env} space, Adler-32 hash-based categorical embeddings (20\,000
buckets), no positional encoding, a 4-layer post-norm transformer, and a
disabled LSTM.

\textbf{Approach.} The training pipeline was run end-to-end with the early
two-stage curriculum (easy: random opponents; medium: 50/50 random and
heuristic opponents) to identify baseline behaviour and failure modes.

\subsection{Validation Infrastructure}

\textbf{Problem.} Without a standardised evaluation protocol, checkpoint
quality could only be assessed by inspecting training-time metrics, which are
noisy and may not reflect true policy strength.

\textbf{Solution.} Implemented a validation pipeline with multiple protocols:

\begin{center}
\small
\begin{tabularx}{\textwidth}{l X}
  \toprule
  \textbf{Protocol} & \textbf{Method} \\
  \midrule
  Smoke        & 10 episodes vs.\ random opponent (quick sanity check). \\
  Fixed paired & 20 frozen teams $\times$ 2 opponents = 40 battles per run. \\
  Mirror       & Both sides use identical team. \\
  \bottomrule
\end{tabularx}
\end{center}

Additionally, 20~frozen Gen\,8 random-battle teams were generated and stored
as JSON manifests (10~swapped pairs for fairness).

\subsection{Vocabulary \& ID Migration}

\textbf{Problem.} Adler-32 hashing with modulo 20\,000 produced an estimated
7.5\,\% collision rate among $\sim$1\,800 unique entity strings.  Collisions
mean different Pok\'{e}mon species (or items or abilities) share the same
embedding, degrading the model's ability to distinguish between entities.

\textbf{Solution.} Replaced hash-based IDs with stable dictionary lookups:

\begin{itemize}
  \item Generated explicit vocabularies from Showdown data files.
  \item Mapped each entity to a unique integer index.
  \item Added support for display-form aliases (e.g.\ ``Shellos (East Sea)''
        $\to$ internal key).
\end{itemize}

Vocabulary sizes: 1\,516~species, 584~items, 315~abilities.
Coverage audit confirmed \textbf{zero unknown entities} across the BDSP
dataset and validation manifests.  This change intentionally broke
compatibility with all previous checkpoints: new training runs are treated as
a new experiment family.

\subsection{Logging \& Metrics Standardisation}

\textbf{Problem.} Metric flow from environment to MLflow had gaps:
``unknown'' opponent metadata was not properly logged, and opponent-specific
metrics were sometimes missing from training runs.

\textbf{Solution.} Traced the full metric path from environment through
\texttt{EnvBridge} to MLflow.  Fixed root causes:
\begin{itemize}
  \item ``Unknown'' opponent metadata now correctly logged as a missing count.
  \item Opponent-specific keys (e.g.\
        \texttt{training/win\_rate\_vs\_random},
        \texttt{training/win\_rate\_vs\_heuristic}) now present in all runs.
  \item Defined a standard contract for high-signal metrics.
\end{itemize}

\subsection{Embedding Audit \& Global Features}

\textbf{Problem.} The observation tensor lacked high-value global context
features that could help the model understand battle dynamics (e.g.\ which
opponent type it was facing, whether it was forced to switch).

\textbf{Solution.} Added to the global token (position~0) without changing
the tensor dimensions:
\begin{itemize}
  \item Opponent label categories (\texttt{opponent\_random},
        \texttt{opponent\_heuristic}, \texttt{opponent\_other}).
  \item Normalised battle turn, force-switch flag, trapped flag.
  \item Normalised available move and switch counts.
  \item Dynamax, Mega Evolution, and Z-Move availability flags.
\end{itemize}

Also discovered and fixed \textbf{broken weather encoding}: the battle state
passed weather as a Python dictionary rather than enum keys, so weather
information was never actually encoded into the observation.

\subsection{Position \& Role Embeddings}

\textbf{Problem.} The transformer received 13~tokens with no positional
information.  Self-attention is permutation-invariant, meaning the model
could not structurally distinguish between the active Pok\'{e}mon and bench
slots, or between own-team and opponent-team tokens.

\textbf{Solution.} Added two sets of learnable embeddings after input
projection:
\begin{itemize}
  \item \textbf{Position embeddings:} 13~vectors, one per token position.
  \item \textbf{Role embeddings:} 5~coarse roles (CLS, our active, our bench,
        opponent active, opponent bench), shared across positions within a role.
\end{itemize}

Old checkpoints are incompatible due to new embedding parameters, so a fresh
training run was required.

\subsection{Compressed Action Space}

\textbf{Problem.} The native \texttt{poke-env} action space has 22~actions
(6~switches + 4~regular moves + 12~gimmick variants).  This large branching
factor slows exploration, and the 12~gimmick variants are rarely all available
simultaneously.

\textbf{Solution.} Introduced a compressed 14-action space:

\begin{center}
\small
\begin{tabularx}{0.7\textwidth}{c l l}
  \toprule
  \textbf{Indices} & \textbf{Category} & \textbf{Mapping} \\
  \midrule
  0--3  & Regular moves & Move 1--4 \\
  4--7  & Gimmick moves & Auto: Mega $\to$ Z $\to$ Dynamax \\
  8--13 & Switches      & Party slot 1--6 \\
  \bottomrule
\end{tabularx}
\end{center}

Gimmick slots automatically resolve to the first legally available gimmick
type, so the agent only needs to learn ``use gimmick'' rather than ``use
specific gimmick variant.''  Action mask changed from $(22,)$ to $(14,)$.
Incompatible with old checkpoints.

\subsection{LSTM Memory Fix}
\label{sec:lstm_memory_fix}

\textbf{Problem.} The LSTM path existed in model code but was non-functional
due to RLlib's connector path not supporting recurrent state propagation in
the production setup.

\textbf{Solution.} Rewrote the LSTM integration:
\begin{itemize}
  \item LSTM now runs even when no incoming recurrent state is provided
        (creates zero initial state).
  \item \texttt{compute\_values()} uses the same forward path as policy logit
        computation, ensuring consistent state handling.
  \item Single-step and sequence batches both pass through the recurrent path
        correctly.
  \item \texttt{get\_initial\_state()} returns unbatched state tensors of shape
        $[H]$ for compatibility with RLlib.
\end{itemize}

End-to-end training validation with LSTM enabled remains a next step.

\subsection{PPO Pipeline Corrections}

\textbf{Problem.} After the architectural improvements, a comprehensive audit
revealed that the PPO training pipeline itself had six issues preventing
effective learning, independent of model architecture.

\textbf{Six fixes applied:}

\begin{center}
\small
\begin{tabularx}{\textwidth}{l r X}
  \toprule
  \textbf{Issue} & \textbf{Before $\to$ After} & \textbf{Impact} \\
  \midrule
  Gradient clipping starvation &
    $\text{clip}=0.5 \to 5.0$ &
    Total grad norm $\sim$65 was clipped to 0.77\,\%, effectively zeroing
    all transformer/embedding gradients. \\
  No per-step action signal &
    $0.0 \to 0.2/0.3$ &
    Matchup and action-quality weights restored, providing per-step credit
    assignment (without them, gradients cancelled over episodes). \\
  Entropy collapse &
    $0.005 \to 0.2$ &
    Entropy bonus was negligible ($\sim$0.02 vs.\ $\sim$5.0 total loss).
    Agent stopped exploring. \\
  Discount horizon too long &
    $\gamma=0.99 \to 0.95$ &
    100-step horizon for $\sim$25-turn battles was too long.  20-step
    horizon better matches episode length. \\
  Value function instability &
    $\text{vf\_clip}=25.0 \to 10.0$ &
    Allowed 3$\times$ the reward range per update, destabilising value
    regression. \\
  LSTM crash on missing state &
    ValueError $\to$ zeros &
    \texttt{\_extract\_state} crashed when state was absent instead of
    returning a zero tensor. \\
  \bottomrule
\end{tabularx}
\end{center}

\begin{keyinsight}
Before diagnosing model architecture, it is essential to verify that gradients
are actually flowing.  In this case, an aggressive gradient clip was silently
starving all learnable layers, and the problem was invisible in standard loss
metrics.
\end{keyinsight}

\subsection{Reward Scaling Approach}
\label{sec:reward_scaling}

\textbf{Problem.} Across the entire project history, explained variance had
been stuck at approximately~0.1 (and sometimes negative).  The value function
was essentially predicting a constant, and value loss showed no downward trend.
This was the single largest blocker to policy improvement.

\textbf{Root cause.} Returns were in the range $[-15, +15]$ with no
normalisation.  RLlib normalises advantages but \emph{not} returns, so the
value function was chasing an unscaled, moving target.  With
\texttt{vf\_clip\_param}$=$10.0, the value head could overshoot by 3$\times$
the reward range per update, preventing convergence.

\textbf{Changes applied:}

\begin{center}
\small
\begin{tabularx}{\textwidth}{l r r X}
  \toprule
  \textbf{Parameter} & \textbf{Before} & \textbf{After} & \textbf{Rationale} \\
  \midrule
  \texttt{reward\_scale} & --- & 0.1 &
    Scales returns from $[-15,+15]$ to $[-1.5,+1.5]$.  Makes value
    regression tractable. \\
  \texttt{vf\_clip\_param} & 10.0 & 3.0 &
    Tighter clipping prevents value head from overshooting. \\
  $\lambda$ (GAE) & 0.92 & 0.88 &
    Reduces GAE target variance. \\
  \texttt{hidden\_dim} & 256 & 512 &
    Wider model gives value head more capacity. \\
  \texttt{num\_heads} & 4 & 8 &
    More heads for finer-grained attention. \\
  \bottomrule
\end{tabularx}
\end{center}

\subsection{Fixed Team \& No-Gimmick Format}

\textbf{Problem.} Random battles introduce high entropy in team composition,
making it difficult to isolate whether the agent is learning battle strategy
or merely memorising matchup patterns.  Additionally, Dynamax and Terastallize
mechanics add complexity that distracts from core battle competency.

\textbf{Solution.} Two changes:
\begin{itemize}
  \item \textbf{Fixed player team:} Added \texttt{player\_team\_path} config
        option.  When set, the agent always uses a specific team (Showdown
        \texttt{.txt} format).  The battle format automatically switches from
        \texttt{gen8randombattle} to \texttt{gen8customgame}.
  \item \textbf{No-gimmick formats:} Created custom Showdown formats
        (\texttt{gen8customgamenogimmicks},
        \texttt{gen8randombattlenogimmicks}) that disable Dynamax, and remove
        Sleep Clause for more natural battle dynamics.
\end{itemize}

With a fixed team, the agent only needs to learn battle strategy, not team
adaptation.  This simplified the learning problem and provided a controlled
setting for validating the training pipeline.

\subsection{Self-Play Infrastructure}

\textbf{Problem.} Self-play opponents were not actually using the training
model's weights, playing with random initialisation instead.

\textbf{Root causes:}
\begin{enumerate}
  \item \textbf{Weights never loaded:} The
        \texttt{selfplay\_weights\_path} was relative
        (\texttt{checkpoints/selfplay\_latest.pt}).  Ray workers run from a
        temporary directory and could not find the file.  The load failure was
        silently swallowed by a bare \texttt{except Exception: pass}.
  \item \textbf{Stale weights:} Weights were only exported every 150K steps.
        The opponent played with outdated checkpoints.
  \item \textbf{Deterministic opponent:} Argmax action selection allowed the
        training agent to exploit deterministic patterns.
\end{enumerate}

\textbf{Fixes:}
\begin{itemize}
  \item Resolved \texttt{selfplay\_weights\_path} to absolute before passing
        to Ray workers.
  \item Replaced \texttt{except: pass} with explicit error logging.
  \item Weight export now happens every training iteration (not just at
        checkpoint saves).
  \item Replaced argmax with temperature sampling ($\tau=0.8$).
  \item Fixed \texttt{opponent\_type=None} to
        \texttt{opponent\_type="self"}, ensuring correct embedding features.
\end{itemize}

After the fix, all workers successfully loaded weights (previously zero
workers).

\begin{keyinsight}
Any file path passed to Ray workers must be resolved to an absolute path
\emph{before} serialisation.  Ray workers execute from a temporary working
directory and cannot find relative paths.  Silent exception handling can mask
critical failures for thousands of training steps.
\end{keyinsight}

\subsection{NaN Bug Workaround}

\textbf{Problem.} After fixing weight loading, the self-play opponent's
fallback rate was 100\,\%, meaning it still played randomly.

\textbf{Root cause.} PyTorch~2.10 has a bug in
\texttt{TransformerEncoderLayer.forward()} when using a float-valued
\texttt{src\_mask}.  The model's learnable attention bias (an
\texttt{nn.Parameter}) triggered this code path, producing NaN logits every
forward pass.  The \texttt{multinomial} sampler then failed, falling back to
random action selection.

\textbf{Fix.} Replaced \texttt{layer(x, src\_mask=bias)} with a manual
norm-first decomposition: LayerNorm $\to$ self-attention $\to$ residual
$\to$ LayerNorm $\to$ FFN $\to$ residual.  This bypasses PyTorch's fused path
and produces correct values.

\textbf{Diagnostic instrumentation added:}
\begin{itemize}
  \item \texttt{\_diag} accumulator in \texttt{SelfPlayPlayer}: weight load
        count, fallback count, action histogram, top probability, entropy,
        valid action count.
  \item Threaded through wrapper $\to$ env bridge $\to$ trainer, logged to
        \texttt{logs/selfplay\_diagnostics.log} and MLflow.
\end{itemize}

\begin{keyinsight}
Do not trust framework-level operations without testing.  A PyTorch internal
bug in the transformer layer was causing NaN outputs, and only explicit
diagnostic logging revealed it. This was discovered because self-play win rate
was basically the same as against random, while it should have been near 50\,\%.
\end{keyinsight}

\subsection{Expanded Validation System}

\textbf{Problem.} The validation system had several limitations: subprocess
per protocol, only argmax actions, missing
\texttt{random\_no\_switch} opponent, no unified benchmark, and no console
output during training.  This made systematic hyperparameter optimisation
difficult.

\textbf{Solution.} Major overhaul of the validation infrastructure:
\begin{itemize}
  \item \textbf{Benchmark protocol:} 3 opponents $\times$ 50 episodes
        (150~total).  Reports per-opponent win rate with Wilson 95\,\%
        confidence intervals, a weighted skill score (random$=$1.0,
        \texttt{random\_no\_switch}$=$1.5, heuristic$=$2.0), and consistency
        (fraction of tiers above 50\,\% win rate).
  \item \textbf{In-process validation:} \texttt{run\_inprocess\_validation()}
        runs against the live algorithm during training, with no subprocess,
        Ray init, or checkpoint restore overhead.
  \item \textbf{Parallel benchmark:} Episodes distributed across $N$~servers
        via \texttt{ThreadPoolExecutor}.  Each environment gets a sequential
        chunk, environments run in parallel.
  \item \textbf{Stochastic evaluation:} \texttt{-{}-explore} flag samples
        from a masked softmax instead of argmax, evaluating the stochastic
        policy.
  \item \textbf{Console output:} Per-opponent win rate bars with confidence
        interval bounds printed at each validation interval.
\end{itemize}

\subsection{Hyperparameter Sweep Methodology}

\textbf{Problem.} With fixed teams performing well, the remaining performance
gap (particularly against heuristic opponents with random teams) required
systematic hyperparameter optimisation rather than manual tuning.

\textbf{Solution.} Implemented an Optuna-based sweep with TPE sampler:
\begin{itemize}
  \item \textbf{Objective:} Benchmark skill score (weighted win rate across
        three opponent tiers).
  \item \textbf{Fixed curriculum:} Single stage with 20\,\% random,
        40\,\% \texttt{random\_no\_switch}, 40\,\% heuristic.  No stage
        promotion, ensuring consistent evaluation across trials.
  \item \textbf{Search space:} 10~parameters covering learning rate
        (log-uniform), entropy coefficient (log-uniform), discount factor,
        GAE $\lambda$, PPO clip, value-function clip, dropout, reward scale
        (categorical), matchup weight, and action-quality weight.
  \item \textbf{Resume mechanism:} SQLite persistence
        (\texttt{logs/hparam\_study.db}) allows interrupting and resuming
        sweeps.
\end{itemize}
